\section*{Phase 3: Data Preparation}

\subsection*{Task 3.1 Selecting Data}

\subsubsection*{Data Rationale Report}

This report outlines the rationale for selecting specific data attributes for the Influenza Outbreak Prediction project. The dataset contains weekly influenza surveillance data for USA and Canada in North America from the GISRS system. Selected data attributes and their importance:
\begin{itemize}
  \item \textbf{Time-based Attributes:} (\texttt{ISO\_YEAR}, \texttt{ISO\_WEEK}, \texttt{ISO\_SDATE}) -- Required for tracking seasonal patterns and outbreak trends.
  \item \textbf{Influenza Cases:} (\texttt{INF\_ALL}, \texttt{INF\_A}, \texttt{INF\_B}) -- Core indicators of outbreak severity and spread.
  \item \textbf{Influenza-Like Illness Activity:} (\texttt{ILI\_ACTIVITY}) -- Proxy measure for estimating undiagnosed cases.
  \item \textbf{Influenza A Subtypes:} (\texttt{AH1}, \texttt{AH3}, \texttt{ANOTSUBTYPED}) -- Helps identify dominant virus strains each season.
\end{itemize}

Below there is a table giving the description and the reason for selection:

\begin{longtable}{|>{\raggedright\arraybackslash}p{3cm}|>{\raggedright\arraybackslash}p{4cm}|>{\raggedright\arraybackslash}p{4cm}|}
\hline
Attribute & Description & Reason for Selection \\
\hline
\endhead
\hline
\texttt{ISO\_YEAR} & The year of the report & Essential for time-series analysis \\
\hline
\texttt{ISO\_WEEK} & The week of the report & Helps in tracking weekly trends \\
\hline
\texttt{ISO\_SDATE} & Start date of the reporting week & Used for time-based feature extraction \\
\hline
\texttt{INF\_ALL} & Total confirmed influenza cases & Primary target variable for prediction \\
\hline
\texttt{INF\_A} & Influenza A cases & Helps in subtype trend analysis \\
\hline
\texttt{INF\_B} & Influenza B cases & Distinguishes between virus types \\
\hline
\texttt{ILI\_ACTIVITY} & Influenza-like illness activity levels & Indicator of outbreak severity \\
\hline
\texttt{AH1}, \texttt{AH3}, \texttt{ANOTSUBTYPED} & Influenza A subtypes & Useful for detailed forecasting \\
\hline
\end{longtable}

Excluded data attributes and rationale:
\begin{itemize}
  \item \textbf{Raw Lab Results:} Excluded since they are redundant with confirmed case counts, and not directly relevant to prediction (\texttt{SPEC\_RECEIVED\_NB}, \texttt{SPEC\_PROCESSED\_NB}).
  \item \textbf{Geographic Identifiers:} Removed as the dataset focuses only on North America (\texttt{WHOREGION}, \texttt{ITZ}, COUNTRY/AREA/TERRITORY, \texttt{COUNTRY\_CODE}).
  \item \textbf{Non-influenza Disease Data:} Outside the scope of this project (\texttt{INF\_NEGATIVE}, \texttt{BVIC}, \texttt{BYAM}, \texttt{BNOTDETERMINED}).
  \item \textbf{Miscellaneous:} (\texttt{ORIGIN\_SOURCE}).
\end{itemize}

\subsection*{Task 3.2 Cleaning Data}

\subsubsection*{Data Cleansing Report}

Identified data issues:
\begin{itemize}
  \item \textbf{Missing Values}: Some records have missing values in \texttt{INF\_ALL}, \texttt{INF\_A}, and \texttt{INF\_B}, \texttt{AH1}, \texttt{AH3}, \texttt{ANOTSUBTYPED}.
  \item \textbf{Outliers}: Unusual spikes in cases that may be due to reporting errors. Looking at all influenza cases trend in North America by plotting the following line plot:
\end{itemize}

\begin{figure}[htbp]
\centering
\includegraphics[width=\textwidth]{05_influenza_trend.png}
\caption{Influenza cases trend in North America.}
\end{figure}

Spikes can be seen in the year 2009, 2014 and 2017 and can be considered outliers depends on how much the influenza cases in those years deviate from the overall trend. To determine this statistically we can use the IQR (Interquartile Range) method. If the influenza cases fall outside the upper or lower boundaries:

Lower boundary: Q1 -- 1.5 * IQR

Upper boundary: Q3 + 1.5 * IQR

Where IQR = Q3 -- Q1

The output is shown as below:

Outliers detected in \texttt{INF\_ALL}: 321

Outliers detected in \texttt{INF\_A}: 347

Outliers detected in \texttt{INF\_B}: 387

After calculating the outliers, we can create a new dataset without them (df\_without\_outliers). After that we can create 3 datasets as following:
\begin{itemize}
  \item Seasonal all influenza cases (seasonal\_inf\_all\_sum)
  \item The mean of all seasonal influenza cases (season\_mean)
  \item Seasonal all influenza cases without outliers (seasonal\_inf\_all\_sum\_without\_outliers)
\end{itemize}

We can plot the bar chart for these 3 datasets as following:

\begin{figure}[htbp]
\centering
\includegraphics[width=0.8\textwidth]{01_seasonal_cases.png}
\caption{Seasonal influenza cases.}
\end{figure}

\begin{figure}[htbp]
\centering
\includegraphics[width=0.8\textwidth]{02_mean_seasonal_cases.png}
\caption{Mean seasonal influenza cases.}
\end{figure}

\begin{figure}[htbp]
\centering
\includegraphics[width=0.8\textwidth]{03_seasonal_no_outliers.png}
\caption{Seasonal influenza cases without outliers.}
\end{figure}

At the end we can do line plot for the following scenarios:
\begin{itemize}
  \item Winter Total Influenza Cases
  \item Mean Influenza Cases (All Seasons)
  \item Mean Influenza Cases Without Outliers (All Seasons)
  \item Winter Total Influenza Cases Without Outliers
\end{itemize}

\begin{figure}[htbp]
\centering
\includegraphics[width=0.85\textwidth]{04_winter_vs_mean.png}
\caption{Winter total cases vs.\ mean seasonal cases.}
\end{figure}

After analyzing the graphs, we can determine that the spikes were real world events:

2009 -- H1N1 Pandemic ("Swine Flu")
\begin{itemize}
  \item Virus: Influenza A(H1N1)pdm09
  \item Overview: A novel H1N1 influenza virus emerged in the spring of 2009, leading to the first influenza pandemic in over 40 years (\textcite{cdc2009}).
  \item Impact: The World Health Organization (WHO) declared the H1N1 flu a pandemic in 2009. That year, the virus caused an estimated 284,400 deaths worldwide (\textcite{mayo_h1n1}).
\end{itemize}

2017 -- Severe Flu Season
\begin{itemize}
  \item Virus: Influenza A(H3N2)
  \item Overview: The 2017-2018 influenza season was notable for its severity, with influenza-like illness (ILI) activity peaking at 7.5\%, the highest percentage since the 2009 pandemic (\textcite{cdc_2017_2018}).
  \item Impact: The season was marked by high hospitalization rates and elevated influenza activity across the United States (\textcite{blanton2017}).
\end{itemize}

Regarding 2014 there were prominent information indicating an outbreak, taking in consideration this information we can conclude that the outliers can be kept for our dataset.
\begin{itemize}
  \item \textbf{Inconsistent Date Formats}: \texttt{ISO\_SDATE} needs to be standardized to pandas datetime object.
  \item \textbf{Removing null/empty values:} NA/NAN rows were dropped.
\end{itemize}

Cleaning strategies:

\begin{longtable}{|>{\raggedright\arraybackslash}p{4cm}|>{\raggedright\arraybackslash}p{8cm}|}
\hline
Issue & Solution Applied \\
\hline
\endhead
\hline
Missing Values & Forward-fill for short gaps, dropping records for excessive missingness \\
\hline
Outliers & Keeping the outliers due to real world outbreaks. \\
\hline
Date Formatting & Converted to datetime format using pandas \\
\hline
\end{longtable}

\subsection*{Task 3.3 Constructing Data}

\subsubsection*{Data Attribute Report}

Newly created attributed are described by the table below:

\begin{longtable}{|>{\raggedright\arraybackslash}p{3cm}|>{\raggedright\arraybackslash}p{5cm}|>{\raggedright\arraybackslash}p{4cm}|}
\hline
New Attribute & Description & Purpose \\
\hline
\endhead
\hline
Month & Extracted from \texttt{ISO\_SDATE} & Seasonality detection \\
\hline
\texttt{Day\_of\_Week} & Extracted from \texttt{ISO\_SDATE} & Identifies weekday trends \\
\hline
\texttt{INF\_ALL\_LAG\_1} & Total cases from the previous week & Short-term trends \\
\hline
\texttt{INF\_ALL\_LAG\_2} & Total cases two weeks prior & Captures disease progression \\
\hline
\texttt{INF\_ALL\_LAG\_4} & Total cases four weeks prior & Captures disease progression \\
\hline
Season & Extracted from \texttt{ISO\_SDATE} & Identifies the season \\
\hline
\end{longtable}

\subsubsection*{Data Generation Report}

Reason for data generation:
\begin{itemize}
  \item To handle missing values in \texttt{ILI\_ACTIVITY} attribute we used forward fill for missing values to propagate the last valid observation forward.
  \item To enhance feature diversity for machine learning models.
  \item Synthetic Forecasting: Created trend-based synthetic values to model future cases for validation.
  \item Season Forecasting: Created season-based values to model season cases for validation.
\end{itemize}

\subsection*{Task 3.4 Integrating Data}

\subsubsection*{Merged Data Set}

Merging process and final dataset composition:
\begin{itemize}
  \item Combined weekly influenza surveillance data with ILI activity reports.
  \item Integrated time-based attributes for seasonal pattern analysis.
  \item Ensured that all records aligned with epidemiological week numbering.
  \item Rows in total \(\sim\)10,000 (varies based on exclusions)
  \item Columns 16 in total (after merging and transformations)
\end{itemize}

\subsection*{Task 3.5 Formatting Data}

\subsubsection*{Final Formatted Data}

The following data adjustment were made:
\begin{itemize}
  \item Standardized date formats (YYYY-MM-DD).
  \item Added column Total\_Influenza\_Cases with data from \texttt{INF\_ALL}, for consistency.
\end{itemize}

The finalized dataset was cleaned and structured dataset was generated in CSV format ready for machine learning model training for further processing.
